% Problem-section file following ../_template.tex.
\section{Weyl--Heisenberg-covariant SICs in every dimension}
\label{sec:problem-18}

\paragraph{Problem.}
Does every finite dimension admit a symmetric informationally complete
measurement that is a single Weyl--Heisenberg orbit?  For an integer $d\geq2$,
let $\omega_d=e^{2\pi i/d}$ and define shift, phase, and displacement operators
on the computational basis by
\begin{equation}
  X_d\lvert j\rangle=\lvert j+1\!\!\pmod d\rangle,
  \qquad
  Z_d\lvert j\rangle=\omega_d^j\lvert j\rangle,
  \qquad
  D_{p,q}=X_d^pZ_d^q,
  \quad (p,q)\in\mathbb Z_d^2.
  \label{eq:p18-displacements}
\end{equation}
Equation~\eqref{eq:p18-displacements} fixes a phase convention that does not
affect the orbit of rank-one projectors.  The question is whether, for every
$d\geq2$, there is a unit vector $\lvert\phi\rangle\in\mathbb C^d$ satisfying
\begin{equation}
  \bigl|\langle\phi\rvert D_{p,q}\lvert\phi\rangle\bigr|^2
  =\frac{1}{d+1}
  \qquad
  \text{for every }(p,q)\in\mathbb Z_d^2\setminus\{(0,0)\}.
  \label{eq:p18-wh-sic-condition}
\end{equation}
If Eq.~\eqref{eq:p18-wh-sic-condition} holds, the $d^2$ projectors in the
Weyl--Heisenberg orbit of $\lvert\phi\rangle$ form a SIC.

\subsection*{Status}

Unsolved

\subsection*{Source}

Renes, Blume-Kohout, Scott, and Caves explicitly conjecture the existence of a
Weyl--Heisenberg-covariant SIC in every finite dimension
\sourcecite{ref:p18-renes}{RBS+04}.

\subsection*{Progress}

\begin{itemize}
  \item Renes, Blume-Kohout, Scott, and Caves stated
  Eq.~\eqref{eq:p18-wh-sic-condition} in every dimension as Conjecture~1 and
  found numerical solutions through $d=45$
  \sourcecite{ref:p18-renes}{RBS+04}.

  \item Appleby, Bengtsson, Flammia, and Goyeneche reported numerical
  Weyl--Heisenberg SICs in every dimension through $d=181$ and in many larger
  dimensions.  These computations establish individual finite cases, not the
  all-dimensional assertion
  \sourcecite{ref:p18-appleby-data}{ABFG19}.

  \item Appleby, Flammia, and Kopp give a construction for all $d>3$
  conditional on the order-one abelian Stark conjecture and a special-value
  identity for the Shintani--Faddeev modular cocycle.  The hypotheses remain
  unproved, so the construction is not unconditional
  \sourcecite{ref:p18-appleby-stark}{AFK25}.
\end{itemize}

\subsection*{References}

\begin{enumerate}[label={},leftmargin=4.8em,labelsep=0.5em]
  \footnotesize
  \item[\textup{[RBS+04]}]\label{ref:p18-renes}
  J. M. Renes, R. Blume-Kohout, A. J. Scott, and C. M. Caves,
  ``Symmetric Informationally Complete Quantum Measurements,''
  \emph{Journal of Mathematical Physics} \textbf{45}, 2171--2180 (2004).
  \href{https://doi.org/10.1063/1.1737053}{doi:10.1063/1.1737053};
  \href{https://arxiv.org/abs/quant-ph/0310075}{arXiv:quant-ph/0310075}.

  \item[\textup{[ABFG19]}]\label{ref:p18-appleby-data}
  M. Appleby, I. Bengtsson, S. Flammia, and D. Goyeneche,
  ``Tight Frames, Hadamard Matrices and Zauner's Conjecture,''
  \emph{Journal of Physics A: Mathematical and Theoretical} \textbf{52},
  295301 (2019).
  \href{https://doi.org/10.1088/1751-8121/ab25ad}{doi:10.1088/1751-8121/ab25ad};
  \href{https://arxiv.org/abs/1903.06721}{arXiv:1903.06721}.

  \item[\textup{[AFK25]}]\label{ref:p18-appleby-stark}
  M. Appleby, S. T. Flammia, and G. S. Kopp,
  ``A Constructive Approach to Zauner's Conjecture via the Stark
  Conjectures,'' arXiv:2501.03970 (2025).
  \href{https://arxiv.org/abs/2501.03970}{arXiv:2501.03970}.
\end{enumerate}

\subsection*{Comment}

The unresolved step is an unconditional construction, or existence proof, for
Eq.~\eqref{eq:p18-wh-sic-condition} in every $d$.  Problem~19 imposes the
additional requirement that the fiducial be Zauner symmetric.

\subsection*{Tag}

Symmetric informationally complete measurements; Quantum measurement theory;
Qudit systems; Representation theory

\subsection*{ID}

\texttt{op\_6ba929179cc40c0a}
